% figures/w01-four-headings.svg — 같은 속도 하나를 두 좌표계로 읽는다.
%
%   bash _tools/tikz2svg.sh figures/src/w01-four-headings.tex figures/w01-four-headings.svg
%
% 네 칸. (1)(2)(3) 은 u 를 1 로 두고 선수각만 바꾼다 — 바뀌는 것은 NED 성분뿐이다.
% (4) 는 선수각을 (1) 과 같게 두고 v 만 더한다 — 선수각이 같은데 진행 방향이 바뀐다.
%
% 수치 (awk 로 계산, 강의 §1-4 의 표와 같다):
%   (1) psi =  0  u=1 v=0    xn_dot=1.0000 yn_dot=0.0000  |V|=1.0000
%   (2) psi = 90  u=1 v=0    xn_dot=0.0000 yn_dot=1.0000  |V|=1.0000
%   (3) psi = 45  u=1 v=0    xn_dot=0.7071 yn_dot=0.7071  |V|=1.0000
%   (4) psi =  0  u=1 v=0.5  xn_dot=1.0000 yn_dot=0.5000  |V|=1.1180
%       beta = atan2(v,u) = 26.565 deg = chi - psi,  chi = atan2(yn_dot, xn_dot)
%
% 화면 규약 : x^n(북) 이 위, y^n(동) 이 오른쪽. 북에서 psi 인 방향의 화면각은 90 - psi.
% 바디 좌표계는 scope 를 그 화면각만큼 돌려서 그린다. 그 안에서 +x 가 선수이고
% **-y 가 우현**이다 (화면각 90-psi-90). 노드는 transform shape 을 쓰지 않으므로
% 글자는 돌지 않는다.
%
% N 과 E 를 그림에 쓰지 않는다 (사용자 지시, 2026-09-12). 축 이름은 x^n, y^n 이고
% 속도 성분은 그 점 미분이다.
\documentclass[border=5pt]{standalone}
% gnc-style.tex — 이 볼트의 TikZ 그림이 공유하는 스타일.
%
% 저널 그림 규약을 스타일 하나로 굳혀 둔다. 그림마다 선 굵기나 화살촉을
% 다시 정하지 않는다 — 그것이 그림이 제각각으로 보이는 원인이다.
%
% 쓰는 법:  % gnc-style.tex — 이 볼트의 TikZ 그림이 공유하는 스타일.
%
% 저널 그림 규약을 스타일 하나로 굳혀 둔다. 그림마다 선 굵기나 화살촉을
% 다시 정하지 않는다 — 그것이 그림이 제각각으로 보이는 원인이다.
%
% 쓰는 법:  % gnc-style.tex — 이 볼트의 TikZ 그림이 공유하는 스타일.
%
% 저널 그림 규약을 스타일 하나로 굳혀 둔다. 그림마다 선 굵기나 화살촉을
% 다시 정하지 않는다 — 그것이 그림이 제각각으로 보이는 원인이다.
%
% 쓰는 법:  \input{gnc-style}  (figures/src/ 안에서)

\usepackage{tikz}
\usepackage{amsmath}
\usepackage{xcolor}
\usetikzlibrary{arrows.meta, positioning, calc, fit, backgrounds, decorations.pathmorphing}

% ---- 색 --------------------------------------------------------------------
% 색은 강조 하나에만 쓴다. 나머지는 검정.
\definecolor{ink}{HTML}{1A1A1A}
\definecolor{accent}{HTML}{7C3AED}
\definecolor{warn}{HTML}{B91C1C}
\definecolor{soft}{HTML}{666666}

% ---- 선과 화살촉 -----------------------------------------------------------
\tikzset{
  % 신호선. 화살촉은 작고 뾰족한 것 하나뿐이다.
  sig/.style   = {ink, line width=0.5pt, -{Stealth[length=4pt, width=3pt]}},
  wire/.style  = {ink, line width=0.5pt},
  % 강조 경로 — 파선으로 구분한다. 흑백 인쇄에서도 살아야 하기 때문이다.
  asig/.style  = {accent, line width=0.5pt, densely dashed,
                  -{Stealth[length=4pt, width=3pt]}},
  awire/.style = {accent, line width=0.5pt, densely dashed},
  % 블록. 흰 바탕, 직각, 검은 테두리.
  blk/.style   = {draw=ink, line width=0.55pt, fill=white,
                  minimum width=17mm, minimum height=8mm, inner sep=2pt, font=\small},
  % 합산점
  sum/.style   = {draw=ink, line width=0.55pt, fill=white, circle,
                  inner sep=0pt, minimum size=4.6mm, font=\scriptsize},
  asum/.style  = {draw=accent, line width=0.55pt, fill=white, circle,
                  inner sep=0pt, minimum size=4.6mm, font=\scriptsize},
  % 분기점
  dot/.style   = {ink, fill=ink, circle, inner sep=0pt, minimum size=1.5mm},
  % 신호 이름 — 선 위에 작게
  sname/.style = {font=\small, inner sep=1.5pt},
  % 그림 안의 짧은 설명. 그림 안의 글은 최소로 둔다
  note/.style  = {font=\scriptsize, soft, inner sep=1.5pt},
  anote/.style = {font=\scriptsize, accent, inner sep=1.5pt},
  % 축
  axis/.style  = {ink, line width=0.5pt},
  tick/.style  = {font=\scriptsize, soft},
}

% 캡션 한 줄 — 그림 아래 가는 선과 함께
\newcommand{\gnccaption}[2]{%
  \node[anchor=north west, text width=#1, align=left, font=\scriptsize, ink]
        at (current bounding box.south west) {\vspace{2pt}#2};}
  (figures/src/ 안에서)

\usepackage{tikz}
\usepackage{amsmath}
\usepackage{xcolor}
\usetikzlibrary{arrows.meta, positioning, calc, fit, backgrounds, decorations.pathmorphing}

% ---- 색 --------------------------------------------------------------------
% 색은 강조 하나에만 쓴다. 나머지는 검정.
\definecolor{ink}{HTML}{1A1A1A}
\definecolor{accent}{HTML}{7C3AED}
\definecolor{warn}{HTML}{B91C1C}
\definecolor{soft}{HTML}{666666}

% ---- 선과 화살촉 -----------------------------------------------------------
\tikzset{
  % 신호선. 화살촉은 작고 뾰족한 것 하나뿐이다.
  sig/.style   = {ink, line width=0.5pt, -{Stealth[length=4pt, width=3pt]}},
  wire/.style  = {ink, line width=0.5pt},
  % 강조 경로 — 파선으로 구분한다. 흑백 인쇄에서도 살아야 하기 때문이다.
  asig/.style  = {accent, line width=0.5pt, densely dashed,
                  -{Stealth[length=4pt, width=3pt]}},
  awire/.style = {accent, line width=0.5pt, densely dashed},
  % 블록. 흰 바탕, 직각, 검은 테두리.
  blk/.style   = {draw=ink, line width=0.55pt, fill=white,
                  minimum width=17mm, minimum height=8mm, inner sep=2pt, font=\small},
  % 합산점
  sum/.style   = {draw=ink, line width=0.55pt, fill=white, circle,
                  inner sep=0pt, minimum size=4.6mm, font=\scriptsize},
  asum/.style  = {draw=accent, line width=0.55pt, fill=white, circle,
                  inner sep=0pt, minimum size=4.6mm, font=\scriptsize},
  % 분기점
  dot/.style   = {ink, fill=ink, circle, inner sep=0pt, minimum size=1.5mm},
  % 신호 이름 — 선 위에 작게
  sname/.style = {font=\small, inner sep=1.5pt},
  % 그림 안의 짧은 설명. 그림 안의 글은 최소로 둔다
  note/.style  = {font=\scriptsize, soft, inner sep=1.5pt},
  anote/.style = {font=\scriptsize, accent, inner sep=1.5pt},
  % 축
  axis/.style  = {ink, line width=0.5pt},
  tick/.style  = {font=\scriptsize, soft},
}

% 캡션 한 줄 — 그림 아래 가는 선과 함께
\newcommand{\gnccaption}[2]{%
  \node[anchor=north west, text width=#1, align=left, font=\scriptsize, ink]
        at (current bounding box.south west) {\vspace{2pt}#2};}
  (figures/src/ 안에서)

\usepackage{tikz}
\usepackage{amsmath}
\usepackage{xcolor}
\usetikzlibrary{arrows.meta, positioning, calc, fit, backgrounds, decorations.pathmorphing}

% ---- 색 --------------------------------------------------------------------
% 색은 강조 하나에만 쓴다. 나머지는 검정.
\definecolor{ink}{HTML}{1A1A1A}
\definecolor{accent}{HTML}{7C3AED}
\definecolor{warn}{HTML}{B91C1C}
\definecolor{soft}{HTML}{666666}

% ---- 선과 화살촉 -----------------------------------------------------------
\tikzset{
  % 신호선. 화살촉은 작고 뾰족한 것 하나뿐이다.
  sig/.style   = {ink, line width=0.5pt, -{Stealth[length=4pt, width=3pt]}},
  wire/.style  = {ink, line width=0.5pt},
  % 강조 경로 — 파선으로 구분한다. 흑백 인쇄에서도 살아야 하기 때문이다.
  asig/.style  = {accent, line width=0.5pt, densely dashed,
                  -{Stealth[length=4pt, width=3pt]}},
  awire/.style = {accent, line width=0.5pt, densely dashed},
  % 블록. 흰 바탕, 직각, 검은 테두리.
  blk/.style   = {draw=ink, line width=0.55pt, fill=white,
                  minimum width=17mm, minimum height=8mm, inner sep=2pt, font=\small},
  % 합산점
  sum/.style   = {draw=ink, line width=0.55pt, fill=white, circle,
                  inner sep=0pt, minimum size=4.6mm, font=\scriptsize},
  asum/.style  = {draw=accent, line width=0.55pt, fill=white, circle,
                  inner sep=0pt, minimum size=4.6mm, font=\scriptsize},
  % 분기점
  dot/.style   = {ink, fill=ink, circle, inner sep=0pt, minimum size=1.5mm},
  % 신호 이름 — 선 위에 작게
  sname/.style = {font=\small, inner sep=1.5pt},
  % 그림 안의 짧은 설명. 그림 안의 글은 최소로 둔다
  note/.style  = {font=\scriptsize, soft, inner sep=1.5pt},
  anote/.style = {font=\scriptsize, accent, inner sep=1.5pt},
  % 축
  axis/.style  = {ink, line width=0.5pt},
  tick/.style  = {font=\scriptsize, soft},
}

% 캡션 한 줄 — 그림 아래 가는 선과 함께
\newcommand{\gnccaption}[2]{%
  \node[anchor=north west, text width=#1, align=left, font=\scriptsize, ink]
        at (current bounding box.south west) {\vspace{2pt}#2};}


\begin{document}
\begin{tikzpicture}[x=1mm, y=1mm]

\def\S{16}           % 1 m/s = 16 mm
\def\P{47}           % 칸 사이 간격

% ---- 칸 하나를 그리는 매크로 -----------------------------------------------
%  #1 x 오프셋   #2 psi [deg]   #3 u   #4 v   #5 xn_dot   #6 yn_dot
\newcommand{\panel}[6]{%
\begin{scope}[shift={(#1,0)}]
  \pgfmathsetmacro{\A}{90-#2}                 % 선수의 화면각
  % NED 축
  \draw[axis, -{Stealth[length=4pt,width=3pt]}] (0,-3) -- (0,27);
  \draw[axis, -{Stealth[length=4pt,width=3pt]}] (-3,0) -- (27,0);
  \node[font=\scriptsize, anchor=south] at (0,27) {$x^n$};
  \node[font=\scriptsize, anchor=west]  at (27,0) {$y^n$};
  % 바디 좌표계 : 선체와 두 축
  \begin{scope}[rotate=\A]
    \draw[soft, line width=0.4pt, fill=black!6]
         (4,0) -- (1.5,1.8) -- (-3.5,1.8) -- (-3.5,-1.8) -- (1.5,-1.8) -- cycle;
    \draw[soft, line width=0.35pt, densely dashed] (4,0) -- (24,0);
    \draw[soft, line width=0.35pt, densely dotted] (0,-1.8) -- (0,-11);
  \end{scope}
  % 속도 벡터 자체 — 굵고 옅게 깔아 u 와 겹쳐도 둘 다 보이게
  \draw[flow, line width=2.4pt, opacity=0.35, -{Stealth[length=6pt,width=5pt]}]
        (0,0) -- ({#6*\S},{#5*\S});
  % u 와 v — 바디에서 잰 것, 선수 끝에 이어 붙인다
  \begin{scope}[rotate=\A]
    \draw[alng, line width=0.9pt, -{Stealth[length=4pt,width=3pt]}] (0,0) -- ({#3*\S},0);
  \end{scope}
\end{scope}}

% ===================== 네 칸 ==============================================
\panel{0}{0}{1}{0}{1}{0}
\panel{\P}{90}{1}{0}{0}{1}
\panel{2*\P}{45}{1}{0}{0.7071}{0.7071}
\panel{3*\P}{0}{1}{0.5}{1}{0.5}

% ---- 칸마다 다른 것 ---------------------------------------------------------
% (1) 선수가 북쪽. u 가 곧 x^n 방향 성분
\node[font=\small\bfseries, anchor=south west] at (-4,31) {(1)\ $\psi = 0^\circ$};
\node[alng, font=\scriptsize, anchor=east] at (-1,9) {$u$};
\node[pathc, font=\scriptsize, anchor=east] at (-1,16) {$\dot x^n$};

% (2) 선수가 동쪽. u 는 그대로인데 x^n 성분이 0
\node[font=\small\bfseries, anchor=south west] at ({\P-4},31) {(2)\ $\psi = 90^\circ$};
\draw[soft, line width=0.45pt] ({\P},10) arc (90:0:10);
\node[soft, font=\scriptsize] at ({\P+8.5},8.5) {$\psi$};
\node[alng, font=\scriptsize, anchor=south] at ({\P+5},0.6) {$u$};   % 호 끝(10,0)을 피한다
\node[pathc, font=\scriptsize, anchor=north] at ({\P+16},-1) {$\dot y^n$};

% (3) 선수가 북동. u 하나가 두 성분으로
\node[font=\small\bfseries, anchor=south west] at ({2*\P-4},31) {(3)\ $\psi = 45^\circ$};
\draw[soft, line width=0.45pt] ({2*\P},14) arc (90:45:14);
\node[soft, font=\scriptsize] at ({2*\P+5},15.5) {$\psi$};
\draw[pathc, line width=0.4pt, densely dashed] ({2*\P+11.314},11.314) -- ({2*\P},11.314);
\draw[pathc, line width=0.4pt, densely dashed] ({2*\P+11.314},11.314) -- ({2*\P+11.314},0);
\node[pathc, font=\scriptsize, anchor=east]  at ({2*\P-1},11.314) {$\dot x^n$};
\node[pathc, font=\scriptsize, anchor=north] at ({2*\P+11.314},-1) {$\dot y^n$};
\node[alng, font=\scriptsize, anchor=north west] at ({2*\P+6},5) {$u$};

% (4) 선수는 (1) 과 같이 북쪽인데 옆으로 밀린다
\node[font=\small\bfseries, anchor=south west] at ({3*\P-4},31)
     {(4)\ $\psi = 0^\circ$, pushed sideways};
%  v 는 u 의 끝에서 우현(동쪽)으로. psi = 0 이면 바디 축과 NED 축이 겹치므로
%  x^n 쪽 투영선은 v 화살표 자체와 같은 선이다 — 따로 긋지 않는다.
\draw[alng, line width=0.9pt, -{Stealth[length=4pt,width=3pt]}]
      ({3*\P},16) -- ({3*\P+8},16);
\node[alng, font=\scriptsize, anchor=south] at ({3*\P+4},16.4) {$v$};
\node[alng, font=\scriptsize, anchor=east]  at ({3*\P-1},9) {$u$};
\draw[pathc, line width=0.4pt, densely dashed] ({3*\P+8},16) -- ({3*\P+8},0);
\node[pathc, font=\scriptsize, anchor=east]  at ({3*\P-1},16) {$\dot x^n$};
\node[pathc, font=\scriptsize, anchor=north] at ({3*\P+8},-1) {$\dot y^n$};
%  크랩각 : 선수(화면 90도) 와 진행 방향(화면 63.43도) 사이
\draw[accent, line width=0.55pt] ({3*\P},12) arc (90:63.435:12);
\node[accent, font=\scriptsize] at ({3*\P+4.2},13.6) {$\beta$};

% ===================== 칸 아래 수치 =========================================
\foreach \i/\uu/\vv/\xd/\yd in {0/1.00/0/1.00/0, 1/1.00/0/0/1.00,
                                2/1.00/0/0.71/0.71, 3/1.00/0.50/1.00/0.50} {
  \node[alng, font=\scriptsize, anchor=west]  at ({\i*\P-4},-15) {$u = \uu$,\ \ $v = \vv$};
  \node[pathc, font=\scriptsize, anchor=west] at ({\i*\P-4},-20) {$\dot x^n = \xd$,\ \ $\dot y^n = \yd$};
}
\node[accent, font=\scriptsize, anchor=west] at ({3*\P-4},-25) {$\beta = \operatorname{atan2}(v,u) = 26.6^\circ$};

% ---- 아래 한 줄 : 색의 뜻 -----------------------------------------------------
\draw[soft, line width=0.3pt] (-4,-30) -- ({3*\P+36},-30);
\node[note, anchor=north west, align=left, text width=180mm] at (-4,-32) {%
  \textcolor{alng}{orange}: $u$, $v$, measured along the vessel's own axes \quad
  \textcolor{flow}{blue}: the velocity itself, one vector \quad
  \textcolor{pathc}{green}: $\dot x^n$, $\dot y^n$, its components along North and East \quad
  \textcolor{accent}{violet}: the crab angle};

\end{tikzpicture}
\end{document}
